\section{三镇合一：边军的重心如何移动}
\label{sec:an-lushan-crisis}

天宝十载，安禄山兼领范阳、平卢、河东三镇。把这件事只写成一个人权势骤增，容易遮住它所依托的地理与制度：范阳面向幽燕及东北边地，平卢的军政资源连着东北方向，河东又关乎山西北部的防务。三个节度使名号集中在同一人手中，意味着原先分置的边军、将校关系、仓廪与转运机会，可以在一个指挥中心附近彼此接通。玄宗朝倚重边将以应付契丹、奚及北方防务的选择，并不必然推出叛乱；但当兵权、供给与任用同时向一处汇集，中央在危机中重新分割或制约这些资源的余地便已缩小。%
\footnote{安禄山兼领三镇见《旧唐书·安禄山传》《新唐书·安禄山传》及《资治通鉴》天宝十载条。三书都足以互证兼领的事实；关于兵额、财赋和安禄山对每一项资源的实际控制，主要是安史以后整理出的叙事，不能把传记中的数目直接当作可核算的常备军或岁入。}

这一段材料需要先分开阅读。《旧唐书》成于后晋，保存了离唐亡较近的纪传层次，却也把败亡后的经验带回天宝朝；《新唐书》在北宋重修，长于把方镇、官制和人物关系编排为可追溯的格局，亦更容易以宋人的制度眼光整理前代。《资治通鉴》把三镇兼领、相权转移和起兵串成紧凑年序，便于看见先后，却以中原朝廷的治乱与决断为轴，不能代替范阳军营、河北州县和边地社会的记录。因而，李林甫死后杨国忠与安禄山的冲突可以说明中枢与边镇的互信迅速下降，不能被缩写成一场私人嫌隙便足以发动战争。%
\footnote{李林甫卒及杨国忠继相，见《旧唐书·李林甫传》《旧唐书·杨国忠传》与《资治通鉴》天宝十一载十一月条。有关两人如何分别影响安禄山的细节，多为事后归责性的宫廷叙事；较稳妥的结论是原有协调关系发生了改变。}

敦煌、吐鲁番的户籍、赋役、仓储和军政文书，使人看到唐代的军费与边防终须落到名册、仓库、差役和行文的层面；墓志又能偶尔钉住官员的任官、迁徙与死亡日期。它们是校正正史整齐叙述的重要旁证，却大多来自河西、西州或保存条件特别优越的精英家庭，不能拿来补成范阳三镇的完整账簿。吐蕃敦煌文书以及突厥、回纥碑铭所保存的边疆语言，也提醒我们：唐史所称的“蕃”“胡”“归附”是朝廷观察中的分类，不是可以直接换算成固定族群或稳定主从关系的标签。安禄山起兵前的东北边地，正因为缺少能够与唐廷记载连续对读的地方档案，尤其不宜用后来的叛乱结局倒推每一个军人、商人或州县居民当时的选择。

\section{范阳南下与两京之间的断裂}

天宝十四载十一月，安禄山自范阳起兵南下。叛军在次月攻陷东都洛阳，取得了关东与河南之间的枢纽，也把唐廷的防线压向潼关。这里的速度本身比任何“盛世突然崩塌”的修辞更值得注目：临时调来的军队、地方守备和朝廷命令，未能在洛阳以北形成足以阻截的合力。可是，洛阳入城的月日、守军的反应和城内秩序，在本纪与列传之间已有差异；史料更少保存普通住民如何避乱、缴纳何种粮饷，或怎样同两方军队打交道。可以确认的是交通与仓储枢纽易手，不能据此把整座城市写成同一种欢迎、服从或抵抗。%
\footnote{起兵与洛阳失陷，见《旧唐书·玄宗本纪下》《旧唐书·安禄山传》《新唐书·安禄山传》及《资治通鉴》天宝十四载十一、十二月条。两《唐书》与《通鉴》构成基本年序；城内社会反应和守军人数的叙述则须保留其来源差异。}

潼关不是一条地图上的细线，而是关东道路、山河地势与关中补给相互收束的关口。洛阳失守后，哥舒翰所部据险防守，唐廷却仍面对收复东都的压力。天宝十五载六月，哥舒翰出关作战而败，燕军得以西进，玄宗随即离开长安入蜀。后世叙事常将这一转折归结为杨国忠逼战，或归咎于哥舒翰一人的进退；《旧唐书·哥舒翰传》与《资治通鉴》的叙述确实保留了将相冲突的线索，但它们不足以让我们精确复原每一道军令的发出时刻。较可把握的是，关中朝廷把一支本可凭险待援的部队置于野战，指挥摩擦、补给压力和政治焦虑在这一决定中同时显露；战败的代价是关中门户洞开。%
\footnote{潼关战败、玄宗入蜀的年序见《旧唐书·肃宗本纪》《旧唐书·哥舒翰传》及《资治通鉴》天宝十五载六月条。关于谁“逼”谁出关，纪传与编年叙事各有归责重点，故只能说中枢与守将间存在严重冲突，不能把复杂的军事决定化约为单一人物的意志。}

长安的失守因此不只是都城换手。皇帝、近臣和部分官署南行，关中居民则面对道路断裂、军队征发和秩序重置；留在城中的人远比逃亡的显要人物难以在史书中被看见。旧、新《唐书》把京师危机纳入玄宗、肃宗两朝的政治叙事，《通鉴》将其安排为政令与军情相继失效的节点。三者所能确证的是制度中心遭到严重破坏，而非一幅可以无缝复原的长安日常图景。城址、墓志和后来可见的迁徙线索可以提示毁坏、流动和复建的不同节奏；它们也不能替我们统计战乱中的死亡或把一处遗存外推为整座京城。

\section{灵武：在行军、诏令与外援之间重立朝廷}

长安失守后，太子李亨在灵武即位，遥尊玄宗为太上皇。至德元载七月的这场即位，首先是一项在军队与官僚网络之间重建名义中心的行动。朔方军及北方官员需要能够发号施令、任命将领并与外部政权交涉的朝廷；远在蜀地的玄宗则难以直接统合北方战场。肃宗朝留下的本纪和诏令传统，把劝进、即位与尊号安排得很有秩序；玄宗是否事先完全同意、各方在灵武的犹疑如何消除，主要仍只能透过胜利者保留下来的文本猜测。写作时应当承认这层不确定，而不能把灵武描绘成毫无裂缝的继承仪式。%
\footnote{灵武即位见《旧唐书·肃宗本纪》《新唐书·肃宗本纪》及《资治通鉴》至德元载七月条。它们足以确认即位、尊玄宗为太上皇及大致程序；劝进者的动机与玄宗预先知情程度，材料并不允许作确定的心理判断。}

灵武朝廷的军事处境，也决定了它不能只靠恢复旧有府兵或命令各州自行勤王。它必须把残余唐军、地方守备、朔方力量与外援重新接到同一条作战线上。张巡在雍丘的守御、其后睢阳的长期围困，显示河南与江淮之间的道路、粮食和城防仍是战争的实际内容；睢阳终于在至德二载十月陷落，不能被忠烈故事磨成只有一方意志的场景。传记中关于饥荒、兵数和守城言辞的笔法极强，较可靠的是围攻持久、守军粮尽、张巡与许远被俘遇害；至于每一项惨烈细节，必须让位于材料的文体和传抄过程。%
\footnote{雍丘与睢阳事见《旧唐书·张巡传》《旧唐书·许远传》《新唐书·张巡传》及《资治通鉴》至德元载、至德二载条。两《唐书》的忠臣传记塑造强烈，宜用以确认守御的过程和政治记忆，不宜据以精算伤亡、存粮或逐字还原言辞。}

至德二载，唐军与回纥骑兵在香积寺击败燕军，随后收复长安。回纥不是被动进入唐朝复国叙事的工具：唐廷需要其机动力量，回纥方面也有婚盟、绢帛、互市与掠获等自身的利益计算。汉文史书对回纥军入关、战利品和市场冲突的记述带有京师居民的焦虑，回纥碑铭及边疆材料又不能逐一补足香积寺的战场细节；两类材料的并读至少能阻止我们把盟军写成无条件的“援兵”。收复长安恢复了象征与行政的中心，却也使对盟友的支付、城内秩序和继续东进的军费立即成为朝廷的难题。%
\footnote{香积寺之战与收复长安见《旧唐书·肃宗本纪》《旧唐书·回纥传》及《资治通鉴》至德二载九月条。唐方所记兵数、战利品和入城秩序，不足以代表回纥方面的全部目标；对回纥的政治行动还须参看其碑铭及其他边疆材料。}

\section{战争没有回到天宝以前}

安禄山死于至德二载正月，其子安庆绪继立；相州邺城的会攻失败后，史思明又杀安庆绪，自立为燕帝；上元二年，史朝义杀史思明而继统。这一连串继承暴力常被写作叛军“自相残杀”，仿佛唐廷只须等待敌人瓦解。实际上，每一次更替都牵动范阳、洛阳、河北前线将领和地方资源的重新站队；唐军同样未能形成稳定的统一指挥。乾元元年，郭子仪等九节度使会攻邺城而溃，正说明把多支军队集中到一处，并不等于能够统一补给、号令和战场判断。%
\footnote{安庆绪、史思明与史朝义的继承，见《旧唐书·安禄山传》《旧唐书·安庆绪传》《旧唐书·史思明传》《旧唐书·史朝义传》及《资治通鉴》至德二载、乾元二年、上元二年条。相州会战见《旧唐书·肃宗本纪》《旧唐书·郭子仪传》及《资治通鉴》乾元元年九月条。双方将领传记相互归咎，兵数与溃败责任不宜照录为定论。}

在这种战争里，财政的重建不是战后附录，而是能否继续作战的条件。乾元元年，刘晏受命整顿盐铁转运，设法把相对稳定的江淮财赋重新输送给中央。旧、新《唐书》的刘晏传和《旧唐书·食货志下》让我们看见盐、漕运与使职网络逐渐成为军费的枢纽；它们也要求我们克制：刘晏在这一年的权限与措施先后并非完全一致，不能把后来盐法的成效提前安到758年。能够明确的是，河北、河南传统税源受战事破坏后，朝廷越来越要依靠江淮的粮运、盐利和专使来供给两京与军队，财政重心由此发生了方向性的移动。%
\footnote{刘晏整顿盐铁转运见《旧唐书·刘晏传》《新唐书·刘晏传》及《旧唐书·食货志下》。志书与传记可说明制度和任用的轨迹，不能据其概括断定同一年内各道的征收比例、盐价或输送总额。}

宝应二年，唐军与回纥军进攻洛阳，史朝义败亡，安史叛乱的主战事至此结束。所谓“结束”只适用于燕政权的主要战场，并不表示天宝旧制重新运转。同年吐蕃军又一度进入长安，代宗出奔；西北边防的空缺说明复都没有抹平战争对军力部署的损耗。对于吐蕃入长安，汉文正史常以京师蒙难为中心，吐蕃方面材料又有日期、行军和内部政治上的缺环。稳妥的写法是承认唐廷的西部防线无力及时合拢，而不是以唐人的屈辱感替代吐蕃行动的全部目的。%
\footnote{史朝义败亡与唐军重入洛阳，见《旧唐书·代宗本纪》《旧唐书·史朝义传》及《资治通鉴》宝应二年正月条；吐蕃入长安见《旧唐书·代宗本纪》《旧唐书·吐蕃传》及《资治通鉴》广德元年十月条。吐蕃军停留日数、撤军原因及所立人选在材料中并有缺环，不能只按唐朝本纪的屈辱叙事处理。}

\section{从战时军镇到可持续的财政与军权}

唐廷为尽快停止河北战事而任用、承认一批旧将，这是一种在资源不足时换取暂时秩序的选择，而不是中央完全不知道其风险。广德二年田承嗣据魏博并收纳安史旧将；大历元年朱希彩杀李怀仙后自领卢龙，朝廷予以承认；大历三年李正己扩展平卢淄青的辖境。三个军镇的地理、军内继承方式和税源并不相同，不能用“藩镇割据”四字抹平。共同之处在于，朝廷收复洛阳后仍须同能够掌握本地士卒、仓储与道路的军事集团谈判，任命文书所写的服从，也不等于州县税收和军队调动已经重新归中央直接支配。%
\footnote{魏博、卢龙与平卢淄青的形成，分别见《旧唐书·田承嗣传》《新唐书·田承嗣传》与《资治通鉴》广德二年条；《旧唐书·朱希彩传》《新唐书·朱希彩传》与《资治通鉴》大历元年条；《旧唐书·李正己传》《新唐书·李正己传》与《资治通鉴》大历三年条。方镇表及中央传记中的州数、兵额多有行政与政治目的，实际赋税和服役结构仍须地方文献、墓志和考古材料补证。}

建中元年杨炎颁行两税法，以夏、秋两次征税来替代已难以维系的租庸调框架。它的前提正是战乱所暴露的现实：户籍散失，授田与人丁相连的旧有征收方式难以贯彻，而中央又必须筹措常备军、转运和京畿防务的费用。两税法把资产、土地与居住地纳入新的征收原则，不应被说成一纸诏令便完成的“近代化”。江南、关中和藩镇辖区的户等编定、钱绢折纳、加派与留截都有不同的执行条件；敦煌等地留存的赋役文书能够帮助我们追问地方实践，却不能直接替代全国账目。%
\footnote{两税法见《旧唐书·杨炎传》《新唐书·食货志五》与《通典》卷七〈食货七·历代盛衰户口〉。诏令和后出的食货志说明改革原则，不等于各地统一税率或即时施行；户等、土地核算及钱绢折纳必须按地区和文书个案分别检验。}

因此，755年的叛乱不能只作为唐由盛转衰的一声断裂来理解。三镇兼领使边军资源的集中成为可能，潼关的失败使都城与中枢同时失去依托，灵武则以军队、诏令与外援重造了一个能够继续战争的朝廷。其后的盐铁转运、河北军镇与两税法，并非同一条直线上的“衰败”或“改革”，而是在兵源、粮道、地方税源和外交代价都已改变的条件下作出的不同应对。正史和《通鉴》使这些决断留下了清晰的年序；文书、碑志与边疆材料则不断提醒我们，年序之外还有被战乱驱散、被赋役重新登记、或根本没有进入史书的人群。唐朝能够收复两京，却不能也没有以同一方式收回所有这些关系。

\subsection{交叉阅读窗口五：长安春望与战时档案的空白}

至德二载（757）陷身长安的杜甫写《春望》：“国破山河在，城春草木深。”这首五言律诗把城破后的视觉停在山河、草木与家书难至之间，文学力量恰在于它不把苦难化成军功的附录。作者是有官僚关系的诗人，所见是被围困都城中的个人处境；它不能统计死亡、粮价或全体居民的立场，更不能替代范阳、洛阳与江淮的经验。%
\footnote{《春望》一般系于至德二载杜甫滞留长安时。诗的系年与作者处境可说明其见证性，不能把“烽火连三月”换算为全国战争持续时间或户口损失。}

因此须把诗与异质材料并置：《旧唐书》《新唐书》的安禄山、肃宗诸传和《资治通鉴》保存三镇、潼关、灵武与收复两京的骨架；《通典》与《唐会要》可追问军费、转运与使职的制度背景；敦煌、吐鲁番军政文书及墓志提供名籍、仓储、任官和迁徙的零星锚点；吐蕃敦煌文书、回纥与突厥碑铭又防止我们只从唐廷看盟友和敌手。它们不能拼成一份连续的战争总账，尤其不能补出范阳基层社会。吕思勉所提出的社会经济长线问题，在这里应问战争怎样改写赋役、流动和家庭，而不是把杜甫的悲哀或正史的治乱判断当成现成答案。
